\section{Two finite networks with a rank saving}
\label{sec:finite-networks}

We construct two fixed linear networks that exchange two banks of values
while restoring every auxiliary value.  When the values are replaced by
arrays, a common representation on all wires at a gate lets that gate act
pointwise without disturbing the representation.  We explain this invariant
before constructing the subspace labels that make changes of representation
cheap.  The first network uses bit values and rational subspaces; the second
uses Gaussian dyadic values and binary subspaces.  The scalar domain specifies
the values combined at a gate, while the subspace labels are used to construct operators
on each wire's address-indexed array.

Put
\[
 h=100,\qquad v=\binom h3=161700,\qquad
 N=v^3=4227952113000000,\qquad m=h^3=1000000.
\]
Let $\mathcal T$ be the set of three-element subsets of $[h]$.
There are two data banks $X_a,Y_a$, indexed by
$a=(a_1,a_2,a_3)\in\mathcal T^3$.
A wire carries one scalar.  A gate is a fixed linear map on a specified
set of wires; untouched wires retain their values.  The vertices on a
wire are its source, the gates that touch it in time order, and its sink;
its edges join consecutive vertices.  We call the name of a wire its
\emph{role}, so that we can distinguish a fixed place in the network
from the value currently carried there.

\subsection{Scalar operations and restored auxiliary values}

There are three stages.  At stage $j$, fix the other two coordinates and
operate on the $v$ values in each bank obtained by varying coordinate $j$.  Thus
there are $v^2$ separate invocations at each stage.
Each such invocation has its own auxiliary wires, distinct from all
other invocations, including those in other stages.  There are
\[
 I=3v^2=78440670000
\]
invocations.  In a forward invocation, write $x_T$ for the data used as
controls and $y_S$ for the data to be updated; $S,T\in\mathcal T$ are
their varying $j$th coordinates.  These are the logical source and
target banks of the invocation.  For each ordered pair $(S,T)$ that is
neighboring under the rule below, there is an auxiliary wire $A_{ST}$.
There are also central wires $C_i$, $i\in[h]$, and in the second
construction a wire $C_\ast$.  We call all these auxiliary wires
\emph{scratch wires}; their initial values may be arbitrary.

The following maps specify the construction.
\begin{itemize}
\item In the bit construction, scalars lie in $\mathbb F_2$, and
$S,T$ are neighbors when $|S\cap T|=1$.
The copy map adds $x_T$ to $A_{ST}$ for each neighbor pair.
The side injection adds $\sum_{T:\,|S\cap T|=1}A_{ST}$ to $y_S$.
The gather adds $\sum_{T\ni i}x_T$ to $C_i$, and the scatter adds
$\sum_{i\in S}C_i$ to $y_S$.
\item In the complex construction, scalars lie in
$\mathbb Z[i,1/2]$, and neighbors are distinct triples with even
intersection.  Copy is unchanged.  Side injection adds
\[
 -\sum_{T\text{ neighboring }S}\frac{|S\cap T|-1}{2}A_{ST}
\]
to $y_S$.  Gather adds $\sum_{T\ni i}x_T$ to $C_i$ and
$\sum_Tx_T$ to $C_\ast$.  Scatter adds
$(\sum_{i\in S}C_i-C_\ast)/2$ to $y_S$.
\end{itemize}
Each update leaves its control wires unchanged.  The first two rows of
the following schedule remove the effect of the old scratch values;
the next four introduce the source values and add their effect to the
target; the last two restore the scratch values.
This cancellation of arbitrary initial scratch values is a linear
instance of the restoration construction for transparent computation in
Buhrman, Cleve, Kouck\'y, Loff and Speelman~\cite[Section~3]{BuhrmanCleveKouckyLoffSpeelman2014}.
The forward schedule is
\begin{equation}
\begin{array}{c|l|l}
\text{row}&\text{operation}&\text{gate grouping}\\ \hline
0&\text{subtract side injection}&\text{one gate per target}\\
1&\text{subtract scatter}&\text{one gate on targets and center}\\
2&\text{copy}&\text{one gate per source}\\
3&\text{gather}&\text{one gate on sources and center}\\
4&\text{scatter}&\text{one gate on targets and center}\\
5&\text{side injection}&\text{one gate per target}\\
6&\text{undo gather}&\text{one gate on sources and center}\\
7&\text{undo copy}&\text{one gate per source}.
\end{array}
\label{eq:motif-schedule}
\end{equation}

\begin{lemma}\label{lem:scalar-motifs}
For arbitrary initial values of its auxiliary wires, a forward
invocation changes $y_S$ to $y_S+x_S$ and leaves every other value
unchanged.  Performing a forward invocation from $X$ to $Y$ at stage 1,
the inverse forward schedule from $Y$ to $X$ at stage 2, and a forward
invocation from $X$ to $Y$ at stage 3 sends
\[
 (X_a,Y_a)\longmapsto(-Y_a,X_a)
\]
and restores every auxiliary value.  In the bit construction this is the
permutation exchanging $X_a$ and $Y_a$.
\end{lemma}

\begin{proof}
Write $\mathcal V$ for copy, $\mathcal G$ for gather,
$\mathcal J$ for side injection and $\mathcal R$ for scatter.
If the initial side and central vectors are $a,c$, the first two rows
subtract $\mathcal J a+\mathcal R c$ from $y$.  Rows 2 and 3 change
them to $a+\mathcal Vx,c+\mathcal Gx$.  Rows 4 and 5 therefore add
$\mathcal R(c+\mathcal Gx)+\mathcal J(a+\mathcal Vx)$.
Rows 6 and 7 restore $c,a$.  The net target change is
$(\mathcal J\mathcal V+\mathcal R\mathcal G)x$, independently of
$a,c$.

In the bit case, the central coefficient from $x_T$ to $y_S$ is
$|S\cap T|\bmod2$.  The side contribution cancels exactly the
intersection-one cases; distinct triples with intersection zero or two
already have coefficient zero, and the self coefficient is one.
In the complex case, the central coefficient is $(|S\cap T|-1)/2$.
The side cancels it for distinct triples with intersection zero or two;
intersection one contributes zero, and the self coefficient is one.
Thus both net maps are $y\gets y+x$.

The inverse schedule with logical source $Y$ and target $X$ effects
$X\gets X-Y$.  Starting from $(x,y)$, the three stages give
$(x,y+x)$, then $(-y,y+x)$, then $(-y,x)$.
Each invocation restores its own auxiliary wires.
\end{proof}

A fixed triple has
\[
 z_{\rm b}=3\binom{97}{2}=13968,\qquad
 z_{\rm c}=\binom{97}{3}+3\cdot97=147731
\]
neighbors in the two constructions, respectively.  The second count
separates intersection zero and intersection two.  Ordered pairs must
be counted: there are $vz$ side wires in each invocation, not $vz/2$.
With $c_{\rm b}=100$ and $c_{\rm c}=101$ central wires per invocation,
the total wire counts are
\begin{align}
 W_{\rm b}&=2N+I(vz_{\rm b}+c_{\rm b})
          =177176569091445000000,\label{eq:bit-wire-count}\\
 W_{\rm c}&=2N+I(vz_{\rm c}+c_{\rm c})
          =1873807244643542670000.\label{eq:complex-wire-count}
\end{align}
All these numbers are fixed, independently of any eventual input length.

\subsection{From scalar values to arrays}

We record the representation principle before choosing the labels.  Let
$R$ be the scalar ring of either network and let $\Omega$ be a finite address
set common to all wires.  Replace each scalar by an array in $R^\Omega$, and
apply each scalar gate pointwise at every address.  Write $\mathsf S_\Omega$
for the resulting network on the vector of wire arrays.

At each source, gate and sink $v$, assign an invertible $R$-linear operator
$D_v$ on $R^\Omega$, called its \emph{frame}.  A frame acts on the values
indexed by addresses within one wire's array and may mix them; a gate instead
combines wire values at each fixed address.  All incoming and outgoing
incidences of one gate use the same $D_v$.  On an edge from $u$ to $v$, apply
$D_vD_u^{-1}$ to that wire's array.  Define the operators on all input and
output roles by
\[
 D_{\rm in}=\bigoplus_w D_{\mathrm{source}(w)},\qquad
 D_{\rm out}=\bigoplus_w D_{\mathrm{sink}(w)}.
\]
For a supplied physical input $f$, define its logical interpretation to be
$g=D_{\rm in}^{-1}f$; this definition requires no physical preprocessing.
At a vertex $v$, the invariant is that the stored array equals $D_v$ applied
to the logical array there.  It holds at the sources by definition.  On an
edge whose logical array is $\varphi$, the edge operator changes $D_u\varphi$
to $D_v\varphi$.  At a gate with logical input arrays $\varphi_j$,
$R$-linearity gives
$\sum_j c_jD_v\varphi_j=D_v(\sum_j c_j\varphi_j)$ for each output combination, with
$c_j\in R$.  Thus the common frame commutes with the pointwise gate, and
the invariant continues to the sinks.  The implemented operator is therefore
\begin{equation}
 D_{\rm out}\,\mathsf S_\Omega\,D_{\rm in}^{-1}.
 \label{eq:common-frame-identity}
\end{equation}
In particular, if the scalar network is a signed role permutation that
sends input role $w$ to output role $\rho(w)$ with scalar sign
$\varepsilon_w$, that route applies
$\varepsilon_wD_{\mathrm{sink}(\rho(w))}
D_{\mathrm{source}(w)}^{-1}$ to its input array.  All roles, including
the auxiliary roles, enter this identity.  Once the scalar routing is fixed,
the endpoint frames determine the array operator on each route; the
intermediate frames can then be chosen to make each edge change inexpensive.

Figure~\ref{fig:frame-compatibility} summarizes this invariant.
\begin{figure}[ht]
\centering
\input{figures/proof-map.tex}
\caption{An edge changes the representation of one role's logical array
$g_w$. At the next gate $G$, all $t$ input arrays have the common frame
$D_v$, so the pointwise linear gate commutes with that frame. The same
$D_v$ acts on each output component.}
\label{fig:frame-compatibility}
\end{figure}

\subsection{Subspaces attached to the gates}

The scalar network already has the required signed exchange. We now
choose its intermediate labels to make the changes of frame inexpensive.
A label is a nondegenerate subspace of a fixed bilinear space. On every
edge one label will contain the other. For such an edge, the
\emph{orthogonal residual} is the orthogonal complement of the smaller label inside
the larger one; its dimension is the absolute change in label dimension.
The saving will come from keeping the total dimension lost on decreasing
edges small relative to $N$.

For the bit network, let $F=\mathbb Q^h$ with bilinear form
\[
 \langle x,y\rangle=x^{\mathsf T}(I-J/9)y,
\]
where $J$ is the all-one matrix.  This form has eigenvalues $1$ and
$1-h/9=-91/9$, so it is nondegenerate.  For a triple $T$ let $t_T$ be
its indicator vector.  Then
\[
 \langle t_S,t_T\rangle=|S\cap T|-1,
 \qquad \langle t_T,t_T\rangle=2.
\]
For the complex network, take instead $F=\mathbb F_2^h$ with the
ordinary dot product; here $t_T\cdot t_T=1$.
Thus every triple line is nondegenerate, and neighboring triple lines
are orthogonal, in the label field belonging to either construction.
The scalar domain need not equal the label field.
These degree-one intersection formulas specialize the construction
described by Alon~\cite[Section~3, remark after Theorem~1.1]{Alon1998}
to the prime $2$, with the ground-set size changed to $h=100$.

The ambient label space is $\mathcal F=F^{\otimes3}$ with its tensor
bilinear form.  It has dimension $m$ and is nondegenerate.  Write
\[
 u_a=t_{a_1}\otimes t_{a_2}\otimes t_{a_3},\qquad U_a=\langle u_a\rangle.
\]
The prescribed terminal labels are
\[
\begin{array}{c|cc}
\text{role}&\text{source label}&\text{sink label}\\ \hline
X_a&U_a&\mathcal F\\
Y_a&0&U_a^\perp\\
\text{each scratch role}&0&\mathcal F.
\end{array}
\]
Write $W$ for the wire count in the construction under consideration.
The source dimensions sum to $N$ and the sink dimensions to $Wm-N$.
Thus the signed dimension changes will sum to $Wm-2N$, because the
contributions at gates cancel. If decreasing edges lose a total
dimension $L$, the sum of absolute changes is $Wm-2N+2L$.
We will obtain $L<N/2$. The rational interface below will require
an additional $N$ source ranks, so this inequality will leave a strict
saving in both constructions.

At stage $j$, fix an invocation.  Put
\[
 A=F^{\otimes(j-1)},\qquad
 P=\langle t_{a_1}\otimes\cdots\otimes t_{a_{j-1}}\rangle,
 \qquad B=P^\perp\subset A.
\]
An empty tensor product is the one-dimensional ground field; thus
$B=0$ at stage 1.  Let
$Q=\langle t_{a_{j+1}}\otimes\cdots\otimes t_{a_3}\rangle
\subset F^{\otimes(3-j)}$, again using the empty tensor product at
stage 3.  Since $P$ is nondegenerate,
$A=B\mathbin{\perp}P$.  
We prescribe the following transition for each data wire in the stage,
where $t$ is the indicator of that wire's $j$th triple:
\[
\begin{array}{c|c|c}
 &\text{incoming label}&\text{outgoing label}\\ \hline
 X&A\otimes\langle t\rangle\otimes Q&(A\otimes F)\otimes Q\\
 Y&B\otimes\langle t\rangle\otimes Q&
 \bigl((B\otimes F)\mathbin{\perp}(P\otimes t^\perp)\bigr)\otimes Q.
\end{array}
\]
The outgoing $Y$ label can also be written
$(P\otimes\langle t\rangle)^\perp_{A\otimes F}\otimes Q$.
These transitions expose one tensor factor at each stage. To see that
they fit together, for $j<3$ write
$Q=\langle t_{a_{j+1}}\rangle\otimes Q'$. The next stage has
\[
 A'=A\otimes F,\qquad
 P'=P\otimes\langle t\rangle,\qquad
 B'=(P')^\perp_{A'}
    =(B\otimes F)\mathbin{\perp}(P\otimes t^\perp).
\]
Thus the prescribed outgoing labels are
$A'\otimes\langle t_{a_{j+1}}\rangle\otimes Q'$ on $X$ and
$B'\otimes\langle t_{a_{j+1}}\rangle\otimes Q'$ on $Y$, exactly
the next incoming labels. At stage 1 they start at $U_a$ on $X_a$
and zero on $Y_a$; at stage 3 they end at $\mathcal F$ on $X_a$
and $U_a^\perp$ on $Y_a$, as prescribed at the terminals.

The gate labels are organized by the decomposition
\[
 A\otimes F=(B\otimes F)\mathbin{\perp}(P\otimes F),
\]
with the common summand $B\otimes F$ and the remaining summand
$P\otimes F$ of dimension $h$. Suppress the one-dimensional factor
$Q$ for the moment. On a side wire joining neighboring physical triples
with indicators $t_X,t_Y$, the common component grows to $B\otimes F$,
while the remaining component passes through
\[
 P\otimes\langle t_X\rangle
 \ \subset\ P\otimes t_Y^\perp
 \ \subset\ P\otimes F.
\]
The first inclusion is possible because $t_X\perp t_Y$.
For the central gates we assign $A\otimes F$ on the physical $X$
bank and $B\otimes F$ on the physical $Y$ bank. A central wire
returns once from the former label to the latter, losing exactly
$P\otimes F$. The gate order below confines all decreases to this
return, while the other edges realize the prescribed stage transitions.

Number the successive times in an invocation by
$r=0,\ldots,7$. At stages 1 and 3, time $r$ uses row $r$ of the scalar
schedule~\eqref{eq:motif-schedule}. At stage 2 it uses the inverse of
row $7-r$, with the logical banks exchanged. In either case the gates
touch the physical wires in the order
\[
 Y\text{-side},\quad Y\text{-center},\quad X\text{-side},\quad
 X\text{-center},\quad Y\text{-center},\quad Y\text{-side},\quad
 X\text{-center},\quad X\text{-side}.
\]
A side wire still joins neighboring physical triples in stage 2 because
the neighbor relation is symmetric. The following table assigns the
common gate labels. It suppresses the tensor factor $Q$; for a gate
indexed by a physical bank triple, $t$ denotes that triple's indicator.

\begin{equation}
\begin{array}{c|l|l}
\text{time}&\text{physical wires}&\text{common gate label}\\ \hline
0&Y,\ \text{its side wires}&B\otimes\langle t\rangle\\
1&Y,\ \text{all center wires}&B\otimes F\\
2&X,\ \text{its side wires}&(B\otimes F)\mathbin{\perp}(P\otimes\langle t\rangle)\\
3&X,\ \text{all center wires}&A\otimes F\\
4&Y,\ \text{all center wires}&B\otimes F\\
5&Y,\ \text{its side wires}&(B\otimes F)\mathbin{\perp}(P\otimes t^\perp)\\
6&X,\ \text{all center wires}&A\otimes F\\
7&X,\ \text{its side wires}&A\otimes F.
\end{array}
\label{eq:motif-labels}
\end{equation}
All labels are tensored by $Q$ and viewed as subspaces of $\mathcal F$.
Every displayed summand is nondegenerate.

All scratch sources have label zero and all scratch sinks have label
$\mathcal F$, as prescribed above. The next lemma verifies the stage
boundaries and all intervening edge changes, including these scratch
edges.

\begin{lemma}\label{lem:motif-residuals}
Every edge joins comparable nondegenerate labels.  The only decreasing
edges are the central-wire edges from time 3 to time 4; their dimension
loss is $h$.  Consequently, if $c$ is the number of central wires per
invocation and $W$ the total number of wires, then
\begin{equation}
 L=Ich,\qquad
 \sum_{e}\bigl|\dim U_{\mathrm{head}(e)}-
                       \dim U_{\mathrm{tail}(e)}\bigr|
 =Wm-2N+2L.                                      \label{eq:rank-budget-labels}
\end{equation}
For the binary label space, every nonzero orthogonal residual of an
edge has an orthonormal basis.
\end{lemma}

\begin{proof}
Write $a=\dim A=h^{j-1}$ and $f=h^{3-j}$.  In the table below,
$\mathrm{in}$ denotes the preceding vertex on that data wire: its source
at stage 1, and its last gate in the preceding stage otherwise.  The
interstage calculation above gives its label.  The table lists the
orthogonal residual and its dimension for every edge except the final
data edges to their sinks, whose labels agree and whose residuals are zero.
Tensor each residual by $Q$, except in the last row.  The symbols $t_X,t_Y$ are
the indicators of the neighboring physical triples joined by a side wire.
\[
\begin{array}{l|l|r}
\text{wire and edge}&\text{orthogonal residual}&\text{dimension}\\ \hline
X:\ \mathrm{in}\to2&B\otimes t_X^\perp&(a-1)(h-1)\\
X:\ 2\to3&P\otimes t_X^\perp&h-1\\
X:\ 3\to6,\ 6\to7&0&0\\
Y:\ \mathrm{in}\to0&0&0\\
Y:\ 0\to1&B\otimes t_Y^\perp&(a-1)(h-1)\\
Y:\ 1\to4&0&0\\
Y:\ 4\to5&P\otimes t_Y^\perp&h-1\\
\mathrm{center}:\ \mathrm{source}\to1&B\otimes F&(a-1)h\\
\mathrm{center}:\ 1\to3&P\otimes F&h\\
\mathrm{center}:\ 3\to4&P\otimes F&h\\
\mathrm{center}:\ 4\to6&P\otimes F&h\\
\mathrm{side}:\ \mathrm{source}\to0&B\otimes\langle t_Y\rangle&a-1\\
\mathrm{side}:\ 0\to2&(B\otimes t_Y^\perp)\mathbin{\perp}(P\otimes\langle t_X\rangle)
 &(a-1)(h-1)+1\\
\mathrm{side}:\ 2\to5&P\otimes\langle t_X,t_Y\rangle^\perp&h-2\\
\mathrm{side}:\ 5\to7&P\otimes\langle t_Y\rangle&1\\
\mathrm{scratch}:\ \mathrm{last}\to\mathrm{sink}
 &(A\otimes F)\otimes Q^\perp&ah(f-1).
\end{array}
\]
The last gate on a central wire is at time 6, and the last gate on a side
wire is at time 7.  Both have label $(A\otimes F)\otimes Q$.  Enlarging it
to $\mathcal F$ at the sink gives the last residual in the table, with
$Q^\perp$ taken in the entire future tensor space.
The side edge $2\to5$ uses the neighbor relation
essentially.  Since $t_X\perp t_Y$, we have
$\langle t_X\rangle\subset t_Y^\perp$.  Hence its two labels satisfy
\[
 (B\otimes F)\mathbin{\perp}(P\otimes\langle t_X\rangle)
 \ \subset\
 (B\otimes F)\mathbin{\perp}(P\otimes t_Y^\perp).
\]
The complement of $\langle t_X\rangle$ inside $t_Y^\perp$ is
$\langle t_X,t_Y\rangle^\perp$, of dimension $h-2$; tensoring by
$P$ and then by $Q$ gives the residual in the table.  The other entries
follow directly by splitting $A=B\mathbin{\perp}P$ and
$F=\langle t\rangle\mathbin{\perp}t^\perp$.
Label dimension is nondecreasing along every edge except center $3\to4$, which removes
$P\otimes F\otimes Q$.  The table includes the interstage edges through
$\mathrm{in}$, and the initial and final data labels match their terminals,
so no other dimension changes occur.

The signed sum of dimension changes telescopes at every gate: each
touched wire has one incoming and one outgoing incidence with the same
gate label.  The sum of source dimensions is $N$, while the sum of
sink dimensions is $Wm-N$.  The signed increase is therefore $Wm-2N$.
There are $Ic$ decreasing edges, each of loss $h$; replacing signed
changes by absolute changes adds $2L$.

Now work over $\mathbb F_2$.  A nondegenerate symmetric bilinear space
is nonalternating precisely when it contains a vector of norm one.
The lines $P,Q$ have such vectors.  If $B\ne0$, a coordinate unit
outside the support of the earlier tensor vector belongs to $B$:
that support has size $3^{j-1}<h^{j-1}$ when $j>1$.
The same argument gives a unit vector in $Q^\perp$ when it is nonzero.
A triple complement, and the complement of two neighboring triple
lines, contain coordinate units outside supports of sizes at most
three and six, respectively.  Full tensor spaces contain coordinate
units as well.  In each nonzero tensor summand, tensoring these chosen
norm-one vectors gives a norm-one vector; placing it in that orthogonal
summand gives one in the whole residual.  Each residual is nondegenerate
by the decompositions above.  Thus every nonzero residual is nonalternating.

For completeness, every nondegenerate nonalternating binary symmetric
space has an orthonormal basis.  Split off unit lines until the
remaining space is alternating.  A nonzero nondegenerate alternating
space splits into planes with bases $a,b$ satisfying
$a\cdot a=b\cdot b=0$, $a\cdot b=1$; this follows by choosing a
nonzero $a$, choosing $b$ with $a\cdot b=1$, and taking the orthogonal
complement of their nondegenerate plane.  Retain one unit vector $w$
from the lines already split off.  A plane orthogonal to $w$ can be
absorbed into that line: the three vectors
\[
 w+a,\qquad w+b,\qquad w+a+b
\]
are independent, pairwise orthogonal, and have norm one.
Use one of these new unit vectors in place of $w$ for the next plane;
the remaining planes are still orthogonal to it.  Repeating this
operation absorbs every alternating plane.
\end{proof}

The loss ratios have the exact values
\[
 \frac{L_{\rm b}}{N}=\frac{100}{539},\qquad
 \frac{L_{\rm c}}{N}=\frac{101}{539}.
\]
In particular, both are smaller than $1/2$.  We now turn the dimension
count into the two interfaces needed later.

\subsection{The rational matrix interface}

For a nondegenerate rational subspace $U\subset\mathcal F$, write
$P_U$ for the projection onto $U$ along $U^\perp$.  Assign $P_U$ to
each gate or sink with label $U$.  Assign the source matrix $-P_{U_a}$
to $X_a$ and zero to every other source.
These are fixed rational $m\times m$ matrices; write $M_v$ for the matrix
at vertex $v$.  Define
$\rho(X_a)=Y_a$, $\rho(Y_a)=X_a$, and let $\rho$ fix every auxiliary
wire.

In the tape application, two address chunks are written as vectors $H,D$
of $m$ fields.  For a matrix $M$, the frame $\Phi_M$ moves the array entry
at $(H,D)$ to $(H+MD,D)$, leaving the other address fields unchanged.
The fields use a radix coprime to the fixed denominators, and addition is
modulo their common range.  Hence $\Phi_M$ is an invertible
$\mathbb F_2$-linear permutation of array entries, and
$\Phi_{M'}\Phi_M=\Phi_{M'+M}$.  Using $D_v=\Phi_{M_v}$
in \eqref{eq:common-frame-identity}, the endpoint matrix difference on each
routed role is exactly the shear applied to that role.  The following identity
makes it the same shear $H\gets H+D$ on every role;
Section~\ref{sec:tape-interchange} turns that shear into a field interchange
and implements an edge change with one smaller interchange per unit of its
rational rank.

\begin{proposition}\label{prop:bit-motif-interface}
The bit network has $W_{\rm b}$ wires, uses only pointwise XOR gates,
and sends its input vector to the permutation $\rho$ of that vector.
Its rational matrices are common to all incidences of a gate and obey
\[
 M_{\mathrm{sink}(\rho(w))}-M_{\mathrm{source}(w)}=I_m
 \quad\text{for every input role }w.
\]
The sum of the rational ranks of all edge differences is
\begin{align}
 s_{\rm b}
 &=\sum_{e}\operatorname{rank}_{\mathbb Q}
      (M_{\mathrm{head}(e)}-M_{\mathrm{tail}(e)})\\
 &=W_{\rm b}m-N+2L_{\rm b}
  =177176569088785861287000000<W_{\rm b}m.
\end{align}
\end{proposition}

\begin{proof}
Lemma~\ref{lem:scalar-motifs} gives the scalar permutation, including
all auxiliary roles.  For $U\subset V$ nondegenerate,
$P_V-P_U$ is the identity on $V\cap U^\perp$ and zero on its
orthogonal complement.  Its rank is $\dim V-\dim U$; reversing the
edge negates the matrix without changing its rank.
The only exceptions are the $N$ negative source projections.
At stage 1, $B=0$, and the first gate on $X_a$ has label exactly $U_a$.
Its first edge matrix is therefore $2P_{U_a}$, of rational rank one,
where the corresponding dimension change was zero.
Lemma~\ref{lem:motif-residuals} consequently gives the stated rank sum.
Its deficit from $W_{\rm b}m$ is $N-2L_{\rm b}>0$.

For an input $X_a$ the endpoint difference is
$P_{U_a^\perp}-(-P_{U_a})=I_m$.  Every other role has zero source
matrix and identity sink matrix.  This proves the interface.
\end{proof}

\subsection{The binary phase interface}

Use the complex scalar network and binary labels.  For a binary vector
$x$, let $\operatorname{wt}(x)$ denote its integer Hamming weight, and
write $[b]\in\{0,1\}$ for the integer representative of $b\in\mathbb F_2$.
For a label $U$, let
\[
 q_U(x)=\operatorname{wt}(P_Ux)\pmod4.
\]
The projection here is over $\mathbb F_2$.
Let
\[
 H=\frac1{\sqrt2}\begin{pmatrix}1&1\\1&-1\end{pmatrix},\quad
 C=H\begin{pmatrix}1&0\\0&i\end{pmatrix}H
   =aI+bX,\qquad a=\frac{1+i}{2},\quad b=\frac{1-i}{2},
\]
where $X$ interchanges the two coordinates.  On functions of
$x\in\mathbb F_2^m$ define
\[
 \mathcal C_U=H^{\otimes m}\operatorname{diag}_x(i^{q_U(x)})H^{\otimes m}.
\]
The use of $H$ in this definition is an identity of matrices.  Each entry
of $\mathcal C_U$ and its inverse is a Gaussian integer divided by $2^m$,
so both preserve arrays over $\mathbb Z[i,1/2]$.  The kernels implemented
below have Gaussian dyadic coefficients.

To apply the scalar network to these functions, replace the scalar on
each wire $w$ by an array
$f_w:\mathbb F_2^m\to\mathbb Z[i,1/2]$.  At each address, a gate
applies its scalar linear map to the values on the wires it touches.
The frame in \eqref{eq:common-frame-identity} is $D_v=\mathcal C_U$ at a
vertex with label $U$.  On an edge from label $U$ to label $V$, apply
$\mathcal C_V\mathcal C_U^{-1}$ to that wire's array.
All frames are diagonalized by the same matrix $H^{\otimes m}$.
For an integer $k\ge1$, an address for $k$ columns is
$(x^{(1)},\ldots,x^{(k)})\in(\mathbb F_2^m)^k$; every frame and edge
operator is then the tensor product of its one-column version over
these $k$ coordinates.  Each scalar gate is still applied once at each
joint address, to the vector of values on the roles it touches.

\begin{proposition}\label{prop:complex-motif-interface}
For each edge of the complex network, write $U,V$ for its tail and
head labels.  Its phase-frame difference
$\mathcal C_V\mathcal C_U^{-1}$ is a product of one translation
kernel $aI+bX_v$, or its inverse, per vector in an orthonormal basis
of the edge residual.  Here $(X_vf)(x)=f(x+v)$.
The sum of these basis lengths is
\[
 s_{\rm c}=W_{\rm c}m-2N+2L_{\rm c}
 =1873807244636671267308000000<W_{\rm c}m.
\]
After source and sink diagonal corrections, and undoing the signed
bank exchange, telescoping these phase frames gives $C^{\otimes m}$
on every role, including auxiliary roles.  The same assertion holds
on any positive number of columns by taking tensor products over columns.
\end{proposition}

\begin{proof}
For binary vectors $z,w$,
\[
 \operatorname{wt}(z+w)
 =\operatorname{wt}(z)+\operatorname{wt}(w)
       -2|\operatorname{supp}(z)\cap\operatorname{supp}(w)|.
\]
If $z\cdot w=0$, the intersection cardinality is even, so weight is
additive modulo four.  If $V=U\mathbin{\perp}E$ and $\mathcal B$
is an orthonormal basis of $E$, it follows that
\[
 q_V(x)-q_U(x)=\sum_{v\in\mathcal B}
                 \operatorname{wt}(v)[v\cdot x]\pmod4.
\]
For a decreasing edge the right side is negated.
Each $v$ has norm one, so $\operatorname{wt}(v)$ is odd.  Its
coefficient in the phase, including the sign for the edge direction,
is therefore some $\varepsilon\in\{1,-1\}$ modulo four.  The function
$i^{\varepsilon[v\cdot x]}$ equals
\[
 \frac{1+i^\varepsilon}{2}
 +\frac{1-i^\varepsilon}{2}(-1)^{v\cdot x}.
\]
The identity $H\operatorname{diag}(1,-1)H=X$, tensored over the
coordinates in $\operatorname{supp}(v)$, sends the diagonal character
$(-1)^{v\cdot x}$ to translation $X_v$.  The displayed phase therefore
becomes $aI+bX_v$ or its inverse.  The nonzero residuals have such
bases by Lemma~\ref{lem:motif-residuals}; for a zero residual use the
empty basis.  The same lemma gives the count $s_{\rm c}$.
The deficit is $2N-2L_{\rm c}>0$.

For an edge $e$, write $\mathcal B_e$ for its residual basis and
$K_{e,v}$ for the forward or inverse translation kernel associated
with $v\in\mathcal B_e$.  Taking tensor products preserves their
product, so on $k$ columns
\[
 \bigl(\mathcal C_V\mathcal C_U^{-1}\bigr)^{\otimes k}
 =\prod_{v\in\mathcal B_e}K_{e,v}^{\otimes k}.
\]
Thus each residual vector contributes one factor acting on all $k$
columns, and there are $s_{\rm c}$ such factors over the network.

Let $\rho$ denote the underlying permutation of the complex roles: it
exchanges $X_a,Y_a$ and fixes each scratch role.
By Lemma~\ref{lem:scalar-motifs}, the pointwise scalar network sends an input array at
$X_a$ to $Y_a$ with sign $+1$, an input at $Y_a$ to $X_a$ with sign
$-1$, and each scratch input to its own output with sign $+1$.
Applying \eqref{eq:common-frame-identity} with $D_v=\mathcal C_{U_v}$,
where $U_v$ is the label at vertex $v$, shows that an input role $w$ reaches
output role $\rho(w)$ with that sign and with operator
$\mathcal C_{U_{\mathrm{sink}(\rho(w))}}
 \mathcal C_{U_{\mathrm{source}(w)}}^{-1}$.
For $k$ columns, use $D_v=\mathcal C_{U_v}^{\otimes k}$ on the joint
address set in the same identity.  The route operator is then the $k$th
tensor power of the displayed address operator, multiplied by its scalar
route sign once.
The zero label has frame $I$, while the full label has frame
$C^{\otimes m}$ because $q_{\mathcal F}(x)=\operatorname{wt}(x)$.
Every pair except $X_a\to Y_a$ has these zero and full labels.
For the remaining pair, write $u=u_a$.  Its binary norm is one and
its weight is $3^3=27$.  Orthogonal weight additivity yields
\begin{align*}
 q_{U_a^\perp}(x)-q_{U_a}(x)
 &=\operatorname{wt}(x)-2\operatorname{wt}(u)[u\cdot x]\\
 &=\operatorname{wt}(x)+2[u\cdot x]\pmod4.
\end{align*}
Since $2[u\cdot x]\equiv2\sum_{j\in\operatorname{supp}(u)}x_j
\pmod4$, the operator on $X_a\to Y_a$ is the tensor of $C^{-1}$ on the
27 coordinates of $\operatorname{supp}(u)$ and $C$ elsewhere.
Put $Z=\operatorname{diag}(1,-1)$.  Direct multiplication gives
\[
 C=iZC^{-1}Z,\qquad C^{-1}=-iZCZ.
\]
For $k$ columns, let $Z_{a,k}$ be the tensor product of $Z$ on
those 27 coordinates in every column and the identity elsewhere.
Apply $Z_{a,k}$ before the network at $X_a$, and apply
$i^{27k}Z_{a,k}$ after the network at $Y_a$; the latter scalar is
applied once to the whole role.  There are $27k$ inverse factors, so
the identity $C=iZC^{-1}Z$ turns each of them into $C$.  The operator
on this pair is therefore $C^{\otimes mk}$ as well.
Only the physical $X$ outputs still carry the minus sign of the
scalar exchange.  Negate those outputs and send output role $\rho(w)$
back to role $w$.  Every role, including every scratch role, now has
the required forward tensor product.
\end{proof}

\subsection{Explicit rational bounds for the two recurrences}

The relative rank deficits are
\[
 \eta_{\rm b}=\frac{W_{\rm b}m-s_{\rm b}}{W_{\rm b}m}
       =\frac{339}{22587335000000},\qquad
 \eta_{\rm c}=\frac{W_{\rm c}m-s_{\rm c}}{W_{\rm c}m}
       =\frac{73}{19906842167500}.
\]
For both recurrences we may use the fixed rational exponent
\begin{equation}
 \tau=\sigma=1-2^{-50}.
 \label{eq:explicit-motif-exponents}
\end{equation}
Both exact deficits displayed above are greater than $20\cdot2^{-50}$.
Since $m<2^{20}$ and $\log 2<1$, we have $\log m<20$, whence
\[
 m^{1-2^{-50}}
 =m\exp(-2^{-50}\log m)
 >m(1-20\cdot2^{-50})>m(1-\eta)
\]
for either $\eta=\eta_{\rm b}$ or $\eta=\eta_{\rm c}$.
Consequently
\[
 \frac{s_{\rm b}}{W_{\rm b}}<m^\tau,
 \qquad \frac{s_{\rm c}}{W_{\rm c}}<m^\sigma,
 \qquad 0<\tau=\sigma<1.
\]
Here $\sigma$ is a rational upper bound for the logarithmic exponent
of the complex recurrence.

